\subsection{The rational Jacobian torsion and its dyadic boundary classes}

\begin{statement}{A.164}{Proved}
{An exact sixteen-element torsion group and four separated boundary images}
For \(J=\operatorname{Jac}(\overline D)\), put
\[
 D_0=[P_0-P_\infty],\qquad D_1=[P_1-P_\infty],\qquad
 D_h=[P_{1/2}-P_\infty].
\tag{A.164.1}
\]
Then
\[
                 J(\Q)_{\rm tors}=\langle D_0,D_1\rangle
                          \simeq(\Z/4)^2,
                \qquad D_h=-D_0-D_1.
\tag{A.164.2}
\]
Every difference of distinct boundary points has exact
order four.  In the Abel--Jacobi map based at \(P_0\),
the four boundary images have the following literal
coordinates in this group:
\[
 \begin{array}{c|cccc}
 P&P_0&P_\infty&P_1&P_{1/2}\\ \hline
 [P-P_0]&0&-D_0&D_1-D_0&-2D_0-D_1.
 \end{array}
\tag{A.164.3}
\]
Their images are pairwise distinct in both
\(J(\Q)/2J(\Q)\) and \(J(\Q_2)/2J(\Q_2)\).
Thus the first finite two-Kummer quotient of the
proper abelian extension distinguishes these four
boundary classes, even after restriction to the
physical place two, once a sheet of \(D^\circ\)
has been selected.  No distinctness assertion under
the larger deck conjugacy of the tripod quotient
is included.
\end{statement}

\begin{proof}
\emph{Divisors and independence of the two order-four generators.}
The totally ramified degree-four map \(T\) gives
\[
 \operatorname{div}(T-c)=4P_c-4P_\infty
                 \quad(c=0,1,\tfrac12),\qquad
 \operatorname{div}(W)=P_0+P_1+P_{1/2}-3P_\infty.
\tag{A.164.4}
\]
These are actual divisors of the displayed rational
functions; the local orders at the four points were
computed in A.150.  They imply \(4D_c=0\) and
the relation in (A.164.2).

There is no degree-two rational function on
\(\overline D_{\overline\Q}\).  To verify the
needed nonhyperellipticity explicitly, use the
three regular differentials of A.150.  Their
complete canonical map on the affine open is
\[
        [dT/W^2:dT/W^3:T\,dT/W^3]=[W:1:T].
\tag{A.164.5}
\]
Its projective image is the smooth plane quartic
\[
                 aW^4=T(T-Z)(T-Z/2)Z.
\tag{A.164.6}
\]
Smoothness follows directly: away from \(W=0\)
the \(W\)-derivative is nonzero; at \(W=0\)
the binary quartic on the right has four simple
roots.  The canonical map is therefore precisely
this embedding of the smooth curve.  A degree-two
map to \(\mathbf P^1\) in characteristic zero
would have an involution acting by minus one on
every regular differential, since its invariant
differentials descend to \(\mathbf P^1\) and are
zero.  The canonical map would then identify its
two generic conjugate points, contradicting
(A.164.5).  This proves the assertion.

Consequently \(2[P-Q]\ne0\) for distinct points
\(P,Q\) among the four boundary points: otherwise
a function with divisor \(2P-2Q\) would give
such a degree-two map.  In particular
\(2D_0\), \(2D_1\), and their sum
\(-2D_h\) are all nonzero.  The two classes
\(2D_0\) and \(2D_1\) are therefore independent
over \(\mathbf F_2\).  A relation
\(rD_0+sD_1=0\) with \(r,s\in\Z/4\)
first has even \(r,s\), by doubling; dividing
these even integers by two and using the same
independence makes both zero modulo four.
Thus the subgroup in (A.164.2) has order sixteen.
The same nonhyperellipticity argument, together
with (A.164.4), gives exact order four for every
nonzero boundary difference.

\emph{Good reduction and a matching upper bound on all rational torsion.}
The constants in (A.160.1) are units at 17 and 31.
The equation (A.164.6) is a smooth proper model
over both local integer rings: four is invertible,
\(a\) is a unit, and the four binary roots remain
distinct.  Its Jacobian therefore has good
reduction at both primes.  The elliptic curves
in the actual rational isogeny
\(J\rightleftarrows E\times\operatorname{Res}_{K/\Q}E_b\)
of A.163 have good reduction there too, since
their short-equation coefficients are nonzero
and their discriminants are units.  The two maps
extend to the good abelian schemes, and their
compositions remain multiplication by two.

Both primes split in \(K=\Q(\sqrt2)\); take
\(\sqrt2=6\) modulo 17 and \(\sqrt2=8\)
modulo 31.  Direct counts on the literal three
elliptic equations give
\[
 \begin{array}{c|ccc|c}
 p&\#E(\mathbf F_p)&\#E_b(\mathbf F_p)
                 &\#E_b^\sigma(\mathbf F_p)&\#J(\mathbf F_p)\\ \hline
 17&20&10&10&2000\\
 31&32&32&32&32768.
 \end{array}
\tag{A.164.7}
\]
The last column follows from the actual isogeny,
not an identification of the finite point groups.
For example, for an isogeny over \(\mathbf F_p\)
with finite etale kernel \(G\), the kernel and
cokernel on rational points are respectively
\(G^{\operatorname{Frob}}\) and
\(G/(\operatorname{Frob}-1)G\).  The latter
description uses the surjectivity of
\(\operatorname{Frob}-1\) on the ambient abelian
variety; its derivative is minus the identity,
so it is a surjective isogeny.  These two finite
kernel groups have the same order.  Thus the
orders of the isogenous point groups agree.
Here the isogeny kernel is killed by two and
\(p\ne2\), so it is etale.  This proves the
order comparison used in (A.164.7).

Reduction injects rational prime-to-\(p\) torsion
at a good prime \(p\).  Indeed the formal
reduction kernel is pro-\(p\), and multiplication
by an integer prime to \(p\) is an isomorphism
on that kernel.  Applying this at 17 and 31
bounds the two-primary rational torsion by
\(\gcd(2000,32768)=16\).  Every odd prime
other than 17 and 31 is excluded by the same
two counts.  Seventeen is excluded by the
reduction at 31, and 31 by the reduction at
17.  Thus the entire rational torsion has order
at most sixteen.  The subgroup already
constructed has that order, proving (A.164.2).

\emph{The dyadic separation.}
Write \(L=\Q_2(\sqrt2)\).  The actual point
functor of the Weil restriction gives
\(A(\Q_2)=E_b(L)\).  First,
\[
              E(\Q_2)[2^\infty]=E[2](\Q_2)
                                   \simeq(\Z/2)^2.
\tag{A.164.8}
\]
There is no point of order four.  In the
duplication formula
\[
                  x([2]P)=
                  \frac{(x(P)^2+n^2)^2}
                       {4x(P)(x(P)^2-n^2)},\qquad n=1146,
\tag{A.164.9}
\]
halving \((0,0)\) requires \(x(P)^2=-n^2\),
which would put \(i\) in \(\Q_2\).  Halving
\((\pm n,0)\) gives respectively
\((x(P)^2\mp2nx(P)-n^2)^2=0\), which would
put \(\sqrt2\) in \(\Q_2\).  Both are
impossible.  A point of higher two-power order
would have a multiple of order four.

For the other factor,
\[
                          E_b(L)[2^\infty]=\{O,(0,0)\}.
\tag{A.164.10}
\]
Indeed \(b=(\sqrt2\epsilon M)^2\cdot573\).
The kernel of the squareclass map
\(\Q_2^\times/\Q_2^{\times2}\to L^\times/L^{\times2}\)
is exactly \(\{1,2\}\): expanding
\((r+s\sqrt2)^2\in\Q_2\) forces \(rs=0\).
Neither 573 nor \(-573\) belongs to this kernel.
Their valuations are zero and their unit residues
modulo eight are five and three, respectively.
Thus \(-b\) is not a square in \(L\), excluding
additional two-torsion.  The duplication formula
for \(E_b\) has numerator \((x^2-b)^2\);
halving \((0,0)\) would make \(b\) a square
in \(L\), also impossible.  This proves
(A.164.10).

The target of \(\Phi\) in (A.163.14) therefore
has local two-primary torsion of exponent two.
For any two-primary torsion point \(R\in J(\Q_2)\),
the exact identity \(\Psi\Phi=[2]\) implies
\([4]R=0\).  In particular \(J(\Q_2)\)
has no point of order eight.  A boundary
difference of exact order four cannot equal
\(2R\) there, since that equality would give
\(R\) exact order eight.  The same argument
over \(\Q\) uses (A.164.2).  This proves
the asserted pairwise separation in the two
quotients by twice the point group.

The finite Kummer maps identify those quotients
with subgroups of the respective \(H^1(J[2])\).
Thus the separation concerns the actual first
two-Kummer quotient in (A.163.1), rather than
only abstract torsion classes.  It is stated
in the chosen cover subgroup; further deck
conjugacy in the tripod extension has not been
used to identify or distinguish its images.

The exact reductions, point counts, divisor
relations, differential map and duplication
factorizations are retained in
\path{code/harmonic_jacobian_torsion_certificate.py}
and its output.
\end{proof}
